% Generated by tools/gen_concepts_values.py -- do not edit.
\subsection*{E00 --- Minimum harness case: one seeded female}
\begin{tabular}{@{}ll@{}}
\textit{Mode} & gate \\
\textit{Seed / days} & 20261117 / 1 \\
\textit{Mosquitoes seeded} & 1 female, 0 male \\
\textit{Closes} & Step 0 \\
\textit{Runs a behaviour test today} & full \\
\end{tabular}

\par\smallskip\noindent\textbf{Purpose.} Checks the scenario harness itself with the smallest possible world: one seeded female adult, one day, and no biology. It answers whether a seeded individual is placed, counted and kept by the chain before any behaviour is tested.
\par\smallskip\noindent\textbf{Setup.} One host-seeking female adult of age zero, seeded at a known position, with no humans, no aquatic habitat and no livestock. The run lasts one day. The scenario is the first step of the harness, so it also checks that the manifest declares exactly one scenario key.
\par\smallskip\noindent\textbf{What passing or failing implies.} If the seeded female is not where it was seeded, or the day-0 count differs from the seeded count, the harness is not reproducing its inputs and no later result can be trusted. A pass says only that seeding and the manifest route are faithful; it says nothing about movement or biology.
\par\smallskip\noindent\textbf{Expected behaviour.} Position, sex and state after one day match the seeded values. These are verifiable by construction, not from a biological band. The manifest declares exactly one scenario key, scenario:E00. This is a gate that closes Step 0 (harness and manifest route), with no tolerance band.
\par\smallskip\noindent\textbf{What the current chain does.} The current chain keeps the single adult in its seeded cell for the whole run and records one successful human bite on day 1 and none afterwards; the bite is drawn against the cell-aggregated human population (about 127,000 people), not against the seeded person
\par\smallskip\noindent\textbf{Checked today by the test suite.}
\begin{itemize}
\item Seeding fidelity: the adult population obeys the conservation identity on every day (adults today = adults yesterday - deaths + emerged, starting from the seeded count), so no seeded individual is dropped or invented; the cohort log and the raster output agree on the number of adults
\item The run completes through the manifest route ('validate\_completeness', resolve, binary) with exit code 0
\item Position pin: the agent remains in the seeded cell
\item GIF baseline pin: the transmission GIF of the one-day run has exactly one frame (a single slide); this is the visualiser state to be improved, not a result
\end{itemize}
\par\noindent\textbf{What the GIF must show.} Today the transmission GIF has one frame per simulated day and no scenario overlay. For E00 the informative version would show the single female marker at its seeded position, with the cell boundary drawn in every frame and a day label. A marker that stays inside its seeded cell confirms the position pin by eye; any jump out of the cell would be visible at once.

\subsection*{E01 --- Person beyond the activation radius}
\begin{tabular}{@{}ll@{}}
\textit{Mode} & gate \\
\textit{Seed / days} & 20261117 / 3 \\
\textit{Mosquitoes seeded} & 1 female, 0 male \\
\textit{People} & 1; distance from the first mosquito: 200 m \\
\textit{Closes} & Step 5 \\
\textit{Runs a behaviour test today} & seeding-only \\
\end{tabular}

\par\smallskip\noindent\textbf{Purpose.} Tests whether a host-seeking female with a person a short distance away, but outside her carbon dioxide activation radius, makes contact at all. It is the first check of the perception cascade and the baseline of the current chain, recorded at Step 1.
\par\smallskip\noindent\textbf{Setup.} One host-seeking female adult and one susceptible person placed just beyond the 100 m carbon dioxide activation radius of the female. There is no aquatic habitat, so the female has no breeding site to oviposit in. The run lasts three days.
\par\smallskip\noindent\textbf{What passing or failing implies.} A pass means the perception cascade respects its activation band: a person outside the activation radius is not contacted, and the chain does not bite people the sensory model should exclude. A fail means the chain still bites at distances the individual geometry rules out, so the transmission numbers would reflect the old cell-scale contact rule instead of individual distance. In the field a female does not find a host she cannot sense.
\par\smallskip\noindent\textbf{Expected behaviour.} Zero successful meals, because distance does not allow contact at this range. No oviposition, since there is no breeding site. The female starves by day three, which is the plan's baseline statement. The target is zero human bites, a strict expected failure until Step 5. The baseline pin records today's behaviour (one bite on day one) as a regression reference, not as an expectation.
\par\smallskip\noindent\textbf{What the current chain does.} The current chain keeps the single adult in its seeded cell for the whole run and records one successful human bite on day 1 and none afterwards; the bite is drawn against the cell-aggregated human population (about 127,000 people), not against the seeded person
\par\smallskip\noindent\textbf{Checked today by the test suite.}
\begin{itemize}
\item Seeding fidelity: the adult population obeys the conservation identity on every day (adults today = adults yesterday - deaths + emerged, starting from the seeded count), so no seeded individual is dropped or invented; the cohort log and the raster output agree on the number of adults
\item The run completes through the manifest route ('validate\_completeness', resolve, binary) with exit code 0
\item Baseline pin: the known-failure facts above are asserted, so any change in the chain's behaviour is flagged and must be acknowledged by updating the pin
\item Target gate (strict expected failure): the plan's target is asserted and is expected to fail until Step 5; when it starts passing the test fails and must be promoted
\item Target for E01: zero successful human bites over the run
\item GIF baseline pin: three simulated days give a three-frame GIF; the target of one frame per phase of the day is a strict expected failure until the visualiser is changed
\end{itemize}
\par\noindent\textbf{What the GIF must show.} Today the transmission GIF has one frame per simulated day and no scenario overlay, so the day-one bite is not visible. For E01 show the female marker, the person marker, and a circle of the 100 m activation radius around the female, redrawn in frames per phase of each night. A correct run shows the female wandering without ever entering the circle and no bite marker at any frame.

\subsection*{E02 --- Night trajectory versus net displacement}
\begin{tabular}{@{}ll@{}}
\textit{Mode} & report \\
\textit{Seed / days} & 20261117 / 3 \\
\textit{Mosquitoes seeded} & 1 female, 0 male \\
\textit{People} & 1; distance from the first mosquito: 2000 m \\
\textit{Closes} & none \\
\textit{Runs a behaviour test today} & seeding-only \\
\end{tabular}

\par\smallskip\noindent\textbf{Purpose.} Reports how far a female travels in one night and how far she ends from her start, when the person is far beyond any single night's reach. It checks the limits of the movement law, not contact, so the arrival of the female at the person is not a pass or fail criterion.
\par\smallskip\noindent\textbf{Setup.} One host-seeking female adult and one susceptible person placed far beyond what the female can cover in one night, with no aquatic habitat. The run lasts three days. The person stays out of reach for the whole run unless the walk happens to drift towards it.
\par\smallskip\noindent\textbf{What passing or failing implies.} The bounds say the movement law cannot produce a trajectory longer than the daily flight budget, nor a net shift beyond the exploration disc, so females cannot be carried across the landscape in one step. If either bound is exceeded, the macro-step or the exploration disc is wrong. A female reaching the person is not a failure, because the direction of the walk decides it.
\par\smallskip\noindent\textbf{Expected behaviour.} Report only. Per night, the trajectory length is at most 590 m ('daily\_move\_capacity\_m', [P], the mean distance travelled per day in the literature). Per night, the net displacement is at most 200 m, which is 2 x 'exploration\_range\_m' (100 m, [EP]). The earlier band of 1,410 to 2,000 m is withdrawn because it cannot be reached under these bounds.
\par\smallskip\noindent\textbf{What the current chain does.} The current chain keeps the single adult in its seeded cell for the whole run and records one successful human bite on day 1 and none afterwards; the bite is drawn against the cell-aggregated human population (about 127,000 people), not against the seeded person
\par\smallskip\noindent\textbf{Checked today by the test suite.}
\begin{itemize}
\item Seeding fidelity: the adult population obeys the conservation identity on every day (adults today = adults yesterday - deaths + emerged, starting from the seeded count), so no seeded individual is dropped or invented; the cohort log and the raster output agree on the number of adults
\item The run completes through the manifest route ('validate\_completeness', resolve, binary) with exit code 0
\item Baseline pin: the known-failure facts above are asserted, so any change in the chain's behaviour is flagged and must be acknowledged by updating the pin
\item Target gate (strict expected failure): the plan's target is asserted and is expected to fail until Step 5; when it starts passing the test fails and must be promoted
\item Target for E02: per-night net displacement bound of 200 m
\end{itemize}
\par\noindent\textbf{What the GIF must show.} Today the transmission GIF has one frame per simulated day and no scenario overlay. For E02 show the female track of each night in its own colour, the start point of the night, a 200 m ring around it, and the far person marker. A correct run shows each night's net shift inside its ring, while the track itself can wander and stays within the daily length budget.

\subsection*{E03 --- Several bites per night in a contact cluster}
\begin{tabular}{@{}ll@{}}
\textit{Mode} & gate \\
\textit{Seed / days} & 20261117 / 3 \\
\textit{Mosquitoes seeded} & 20 female, 0 male \\
\textit{People} & 5; distance from the first mosquito: 141 m, 424 m, 707 m, 728 m, 728 m \\
\textit{Closes} & Step 3 \\
\textit{Runs a behaviour test today} & seeding-only \\
\end{tabular}

\par\smallskip\noindent\textbf{Purpose.} Checks that one mosquito can make several blood-feeding attempts in one night and that each attempt is logged as its own event, not folded into one cell-level count. It is the first test of the multi-bite ledger.
\par\smallskip\noindent\textbf{Setup.} A cluster of host-seeking females with a few susceptible people placed among them, so that contacts occur in several places during the run, in an area with no aquatic habitat. The design freezes the probing success rate so that the expected count is a closed-form product. The JSON does not yet carry that override (its overrides are empty); the field enters with Step 3.
\par\smallskip\noindent\textbf{What passing or failing implies.} The ledger must count bites, not cell aggregates. If events per night equal the number of populated cells rather than the number of bites, transmission is understated in dense clusters and the per-person bookkeeping breaks. A pass confirms the multi-bite rule; a fail means the one-bite-per-night simplification is still present.
\par\smallskip\noindent\textbf{Expected behaviour.} Gate. The number of meals equals the number of females times the probing success rate times the probability of contact, computed by hand before the code is written. The ledger records one event per bite attempt, so a night shows several events, not one per cell.
\par\smallskip\noindent\textbf{What the current chain does.} Not meaningful yet: depends on Step 3 (blood-meal physiology and multi-bite pickup).
\par\smallskip\noindent\textbf{Checked today by the test suite.}
\begin{itemize}
\item Seeding fidelity: the adult population obeys the conservation identity on every day (adults today = adults yesterday - deaths + emerged, starting from the seeded count), so no seeded individual is dropped or invented; the cohort log and the raster output agree on the number of adults
\item The run completes through the manifest route ('validate\_completeness', resolve, binary) with exit code 0
\end{itemize}
\par\noindent\textbf{What the GIF must show.} One frame per phase of each night. Each female is a marker and each bite is a flash on the person it hits, with a running count of bites per person for the night. A correct run shows several flashes in one night from the same female, and the flashes land on different people as the females move.

\subsection*{E04 --- Partial meal does not start the cycle}
\begin{tabular}{@{}ll@{}}
\textit{Mode} & gate \\
\textit{Seed / days} & 20261117 / 3 \\
\textit{Mosquitoes seeded} & 1 female, 0 male \\
\textit{People} & 1; distance from the first mosquito: 200 m \\
\textit{Closes} & Step 3 \\
\textit{Runs a behaviour test today} & seeding-only \\
\end{tabular}

\par\smallskip\noindent\textbf{Purpose.} Checks that a female whose blood meal is below target does not start the gonotrophic cycle: she keeps host-seeking that night and nothing is laid.
\par\smallskip\noindent\textbf{Setup.} One host-seeking female and one susceptible person placed next to each other, so a meal is possible. The intent is that the female starts with a partial blood volume below the target. The JSON does not yet carry a blood-volume field, so the partial volume is seeded when Step 3 adds it.
\par\smallskip\noindent\textbf{What passing or failing implies.} Without this rule a single small feed would start the cycle, overstating egg production and the turnover of biting females. If the cycle advances after a partial meal, the physiology is not using the target volume. Real females need a sufficient meal before they mature eggs.
\par\smallskip\noindent\textbf{Expected behaviour.} Gate. After a partial meal the female remains host-seeking that night and the cycle does not start. The threshold is 'blood\_meal\_target\_ul' = 0.5 uL [EP].
\par\smallskip\noindent\textbf{What the current chain does.} Not meaningful yet: depends on Step 3 (blood-meal volume and the cycle clock).
\par\smallskip\noindent\textbf{Checked today by the test suite.}
\begin{itemize}
\item Seeding fidelity: the adult population obeys the conservation identity on every day (adults today = adults yesterday - deaths + emerged, starting from the seeded count), so no seeded individual is dropped or invented; the cohort log and the raster output agree on the number of adults
\item The run completes through the manifest route ('validate\_completeness', resolve, binary) with exit code 0
\end{itemize}
\par\noindent\textbf{What the GIF must show.} One frame per phase. The female is drawn in host-seeking colour, with a blood volume bar beside her that grows with each meal and a cycle timer beside it. A correct run keeps the female in host-seeking colour after the partial meal and the cycle timer stays at zero.

\subsection*{E05 --- Egg batch proportional to accumulated blood}
\begin{tabular}{@{}ll@{}}
\textit{Mode} & gate \\
\textit{Seed / days} & 20261117 / 3 \\
\textit{Mosquitoes seeded} & 2 female, 0 male \\
\textit{People} & 1; distance from the first mosquito: 200 m \\
\textit{Closes} & Step 3 \\
\textit{Runs a behaviour test today} & seeding-only \\
\end{tabular}

\par\smallskip\noindent\textbf{Purpose.} Checks that the egg batch scales with the accumulated blood volume a female actually took, not with a fixed binomial batch, by comparing a partially fed female with a fully fed one.
\par\smallskip\noindent\textbf{Setup.} Two host-seeking females next to the same susceptible person. One is meant to start with a partial blood volume and the other with a full meal. The JSON does not yet seed volumes, so the two females start identical until Step 3 adds the volume field.
\par\smallskip\noindent\textbf{What passing or failing implies.} Egg output must follow the blood volume the female took, which is the link the calibration rests on. If both females lay the same batch, the model loses the feeding-fecundity link that the literature measures, and the egg count no longer responds to host access.
\par\smallskip\noindent\textbf{Expected behaviour.} Gate. The egg count equals the blood volume times 223 eggs per uL [EP], rounded. This is calibrated so that a 0.5 uL meal gives about 112 eggs, close to the 111.67 plus or minus 5.36 measured in coluzzii [P]. The batch is proportional to the accumulated volume, so the fully fed female lays more than the partial one.
\par\smallskip\noindent\textbf{What the current chain does.} Not meaningful yet: depends on Step 3 (egg batch proportional to blood volume).
\par\smallskip\noindent\textbf{Checked today by the test suite.}
\begin{itemize}
\item Seeding fidelity: the adult population obeys the conservation identity on every day (adults today = adults yesterday - deaths + emerged, starting from the seeded count), so no seeded individual is dropped or invented; the cohort log and the raster output agree on the number of adults
\item The run completes through the manifest route ('validate\_completeness', resolve, binary) with exit code 0
\end{itemize}
\par\noindent\textbf{What the GIF must show.} Two females side by side, each with a blood volume bar and an egg count shown at the end of the run. A correct run shows the two bars in the same ratio as the two egg counts.

\subsection*{E06 --- First cycle needs two meals}
\begin{tabular}{@{}ll@{}}
\textit{Mode} & gate \\
\textit{Seed / days} & 20261117 / 3 \\
\textit{Mosquitoes seeded} & 1 female, 0 male \\
\textit{People} & 1; distance from the first mosquito: 200 m \\
\textit{Closes} & Step 3 \\
\textit{Runs a behaviour test today} & seeding-only \\
\end{tabular}

\par\smallskip\noindent\textbf{Purpose.} Checks that a teneral female, which has to complete a first gonotrophic cycle, needs two blood meals, because the first-cycle target is above the usual meal target.
\par\smallskip\noindent\textbf{Setup.} One teneral female adult, freshly emerged, next to one susceptible person. The JSON seeds the teneral state; the two-meal rule comes from the first-cycle and usual meal targets that Step 3 introduces.
\par\smallskip\noindent\textbf{What passing or failing implies.} If one meal completes the first cycle, the model underestimates the time to first oviposition and therefore the pace at which young females transmit. Requiring two meals keeps the young-female pace in line with the physiology the plan adopts.
\par\smallskip\noindent\textbf{Expected behaviour.} Gate. One meal is not enough. The female needs two, because 'first\_cycle\_target\_ul' = 0.8 uL [EP] is above 'blood\_meal\_target\_ul' = 0.5 uL [EP]. The first cycle closes only once the first-cycle target is reached.
\par\smallskip\noindent\textbf{What the current chain does.} Not meaningful yet: depends on Step 3 (first-cycle target of two meals).
\par\smallskip\noindent\textbf{Checked today by the test suite.}
\begin{itemize}
\item Seeding fidelity: the adult population obeys the conservation identity on every day (adults today = adults yesterday - deaths + emerged, starting from the seeded count), so no seeded individual is dropped or invented; the cohort log and the raster output agree on the number of adults
\item The run completes through the manifest route ('validate\_completeness', resolve, binary) with exit code 0
\end{itemize}
\par\noindent\textbf{What the GIF must show.} One frame per phase, with the female's blood volume bar and a marker for the first-cycle target. A correct run shows the bar crossing the usual target after the first meal but not the first-cycle target, and the cycle completes only after the second meal.

\subsection*{E07 --- Breeding site beats the host}
\begin{tabular}{@{}ll@{}}
\textit{Mode} & gate \\
\textit{Seed / days} & 20261117 / 3 \\
\textit{Mosquitoes seeded} & 1 female, 0 male \\
\textit{Aquatic cohorts} & 20 PUPA at (5100, 5200) m \\
\textit{People} & 1; distance from the first mosquito: 2000 m \\
\textit{Closes} & Step 7 and Step 9 \\
\textit{Runs a behaviour test today} & seeding-only \\
\end{tabular}

\par\smallskip\noindent\textbf{Purpose.} Two assertions in one scenario. Assertion A: an emerging adult appears at the position of its mother cohort, not at the centre of the cell. Assertion B: a gravid female chooses the breeding habitat, not the host, when both are on offer. Assertion B needs a gravid female, which the JSON does not seed yet.
\par\smallskip\noindent\textbf{Setup.} One host-seeking female, a nearly emerged pupal cohort with high development a short distance from her, and one person far beyond her reach. The layout puts the breeding site nearer than the host, so the choice is between them.
\par\smallskip\noindent\textbf{What passing or failing implies.} For A, emergence at the cell centre would leak the cell scale into the individual and lose the mother's location. For B, a female that goes to the person ignores the habitat signal, and biting is then driven by proximity rather than by the oviposition choice. Either failure would misplace the transmission the model reports.
\par\smallskip\noindent\textbf{Expected behaviour.} Gate with two assertions. A, closing Step 7: the emergence location equals the mother position, not the cell centre. B, closing Step 9: the oviposition site is the habitat site, not the host. Oviposition in natal habitat is not modelled (section 16.1), so B is about choosing habitat over host, not about the natal site.
\par\smallskip\noindent\textbf{What the current chain does.} Not meaningful yet: depends on Step 7 (emergence at the mother position) and Step 9 (oviposition site choice).
\par\smallskip\noindent\textbf{Checked today by the test suite.}
\begin{itemize}
\item Seeding fidelity: the adult population obeys the conservation identity on every day (adults today = adults yesterday - deaths + emerged, starting from the seeded count), so no seeded individual is dropped or invented; the cohort log and the raster output agree on the number of adults
\item The run completes through the manifest route ('validate\_completeness', resolve, binary) with exit code 0
\item Seeded pupae emerge: between one and the seeded number of adults emerge from the pupae seeded at the breeding site outside the candidate cells, so aquatic seeding at a metric position is live
\end{itemize}
\par\noindent\textbf{What the GIF must show.} One frame per phase. The pupal cohort is a marker with its emergence point, the female has a track, and the person is a far marker. For A the new adult should appear on the pupal marker, not at the cell centre. For B the female's track should end on the habitat marker rather than on the person.

\subsection*{E08 --- Two females, two persons, distinct identities}
\begin{tabular}{@{}ll@{}}
\textit{Mode} & gate \\
\textit{Seed / days} & 20261117 / 3 \\
\textit{Mosquitoes seeded} & 2 female, 0 male \\
\textit{People} & 2; distance from the first mosquito: 60 m, 7114 m \\
\textit{Closes} & Step 2 \\
\textit{Runs a behaviour test today} & seeding-only \\
\end{tabular}

\par\smallskip\noindent\textbf{Purpose.} Checks that each person has a unique, stable identifier and that each female is linked to the person she is matched with. The plan closes it at Step 2 with identifier checks: uniqueness, stability and the number of persons per cell.
\par\smallskip\noindent\textbf{Setup.} Two host-seeking females in two cells far apart from each other, each with a susceptible person well inside her activation radius. The pairs are separated so that no female is close to the other female's person.
\par\smallskip\noindent\textbf{What passing or failing implies.} If person identifiers collide or change between runs, the bite ledger cannot say who was bitten and the individual-level transmission cannot be traced. A pass means each person keeps one identity; it does not by itself prove the choice rule of the cascade.
\par\smallskip\noindent\textbf{Expected behaviour.} Gate at Step 2, verifiable by construction: each person has a unique human identifier that is deterministic from the seed and stable across runs (the earlier cell-times-slots formula was dropped by decision 2a), and the number of persons per cell matches the population. From Step 5 each female bites the person the weighted roulette picks among those inside her perception radius, with distance as the criterion of the stretch, and there is no winning cell.
\par\smallskip\noindent\textbf{What the current chain does.} Not meaningful yet: depends on Step 2 (individual humans with stable identifiers).
\par\smallskip\noindent\textbf{Checked today by the test suite.}
\begin{itemize}
\item Seeding fidelity: the adult population obeys the conservation identity on every day (adults today = adults yesterday - deaths + emerged, starting from the seeded count), so no seeded individual is dropped or invented; the cohort log and the raster output agree on the number of adults
\item The run completes through the manifest route ('validate\_completeness', resolve, binary) with exit code 0
\end{itemize}
\par\noindent\textbf{What the GIF must show.} One frame per phase, with each person drawn in a colour fixed by identifier, so that a person keeps the same colour across frames. A correct run shows each person keeping its colour and each bite landing on exactly one of them.

\subsection*{E09 --- Infective female, several persons}
\begin{tabular}{@{}ll@{}}
\textit{Mode} & gate \\
\textit{Seed / days} & 20261117 / 3 \\
\textit{Mosquitoes seeded} & 1 female, 0 male \\
\textit{People} & 3; distance from the first mosquito: 40 m, 60 m, 64 m \\
\textit{Closes} & Step 2 \\
\textit{Runs a behaviour test today} & seeding-only \\
\end{tabular}

\par\smallskip\noindent\textbf{Purpose.} Checks the one-way rule of transmission at the bite: an infective female gives several infectious bites across the run when several persons are around, the infection can reach more than one person across the set of runs, and each infectious bite infects at most the person it bites.
\par\smallskip\noindent\textbf{Setup.} One infective host-seeking female with three susceptible persons placed around her at different short distances, run for three days. The infective flag is set on the mosquito and all persons start susceptible.
\par\smallskip\noindent\textbf{What passing or failing implies.} If one bite infects anyone other than the person bitten, the model creates transmission that no contact produced and inflates incidence. If the infection never reaches a second person across the runs, the spread between persons is missing and the model cannot show onward transmission.
\par\smallskip\noindent\textbf{Expected behaviour.} Gate at Step 2. The run contains several infectious bites; the probability that more than one person is infected is above zero; each infectious bite infects at most the person bitten. The check is structural, not a fixed band.
\par\smallskip\noindent\textbf{What the current chain does.} Not meaningful yet: depends on Step 2 (individual humans and the per-bite ledger).
\par\smallskip\noindent\textbf{Checked today by the test suite.}
\begin{itemize}
\item Seeding fidelity: the adult population obeys the conservation identity on every day (adults today = adults yesterday - deaths + emerged, starting from the seeded count), so no seeded individual is dropped or invented; the cohort log and the raster output agree on the number of adults
\item The run completes through the manifest route ('validate\_completeness', resolve, binary) with exit code 0
\end{itemize}
\par\noindent\textbf{What the GIF must show.} One frame per phase. The infective female is a red marker, the persons are grey until infected, and each infectious bite turns its person red. A correct run shows red bite flashes that change only the bitten person, never a neighbour.

\subsection*{E10 --- Sensitivity to the three perception radii}
\begin{tabular}{@{}ll@{}}
\textit{Mode} & report \\
\textit{Seed / days} & 20261117 / 3 \\
\textit{Mosquitoes seeded} & 1 female, 0 male \\
\textit{People} & 1; distance from the first mosquito: 150 m \\
\textit{Variants} & activacion\_co2\_50\_100\_200\_m, olfativa\_10\_30\_70\_m, visual\_0.25\_0.72\_m \\
\textit{Overrides} & co2\_activation\_radius\_m, olfactory\_radius\_m, visual\_radius\_m \\
\textit{Closes} & none \\
\textit{Runs a behaviour test today} & seeding-only \\
\end{tabular}

\par\smallskip\noindent\textbf{Purpose.} Quantifies how much each stage of the sensory cascade (carbon dioxide activation, olfactory tracking, visual approach) moves the outcome, by sweeping one radius at a time. It is a report: it measures sensitivity and does not gate.
\par\smallskip\noindent\textbf{Setup.} One host-seeking female and one susceptible person a short distance away, run three days, once per variant. The three variants live in one file, as entries in its variants list, and each varies one radius with the others at their defaults. Variant one sweeps the carbon dioxide activation radius over 50, 100 and 200 m. Variant two sweeps the olfactory radius over 10, 30 and 70 m. Variant three sweeps the visual radius over 0.25 and 0.72 m. The defaults are the verified values.
\par\smallskip\noindent\textbf{What passing or failing implies.} The result shows which radius the outcome depends on. If the outcome barely moves across a sweep, that stage is doing little work and the cascade is poorly identified there; if it moves a lot, the default needs evidence, which the plan marks as [EP] where there is none. This tells the team which radius to calibrate first.
\par\smallskip\noindent\textbf{Expected behaviour.} Report only, with no pass or fail. For each variant the observed contacts and outcome are written next to the default run. The plan gives no band here, because the sweep is itself the result.
\par\smallskip\noindent\textbf{What the current chain does.} Not meaningful yet: depends on Step 5 (the cascade). Today the overrides are not applied: the engine rejects non-empty overrides until Step 5, so the harness runs each variant without them, with the same seed, and records the overrides as not applied. The three variants therefore reduce to the same run and the report cannot show sensitivity.
\par\smallskip\noindent\textbf{Checked today by the test suite.}
\begin{itemize}
\item Seeding fidelity: the adult population obeys the conservation identity on every day (adults today = adults yesterday - deaths + emerged, starting from the seeded count), so no seeded individual is dropped or invented; the cohort log and the raster output agree on the number of adults
\item The run completes through the manifest route ('validate\_completeness', resolve, binary) with exit code 0
\end{itemize}
\par\noindent\textbf{What the GIF must show.} One GIF per variant. Each shows the female track, the person, and the radius circle drawn at the swept value. Once the overrides apply, the three GIFs should differ in circle size and in when contact happens; if they look the same, the sweep is not reaching the engine.

\subsection*{E11 --- Postprandial window after a full meal}
\begin{tabular}{@{}ll@{}}
\textit{Mode} & gate \\
\textit{Seed / days} & 20261117 / 3 \\
\textit{Mosquitoes seeded} & 1 female, 0 male \\
\textit{People} & 1; distance from the first mosquito: 100 m \\
\textit{Closes} & Step 3 \\
\textit{Runs a behaviour test today} & seeding-only \\
\end{tabular}

\par\smallskip\noindent\textbf{Purpose.} Checks that a female that has completed a meal cannot try again before the postprandial window closes, and records how the return to host-seeking rises afterwards. The window is a gate; the reopening curve is reported only.
\par\smallskip\noindent\textbf{Setup.} One host-seeking female, fed on day zero, in the same cell as a susceptible person, so that she would try to feed again at once if no window existed. The JSON does not yet seed a fed state; the fed state comes from the Step 3 physiology.
\par\smallskip\noindent\textbf{What passing or failing implies.} If the female bites again immediately, the blood-feeding rate is too high and transmission is inflated. The reopening curve describes how fast a fed female resumes host-seeking. It is an estimate without a direct source, so it is reported with its band, not gated.
\par\smallskip\noindent\textbf{Expected behaviour.} Gate. The first attempt after the target is reached is not before 12 hours, given 'blood\_meal\_refractory\_days' = 0.5 day [EP], and falls in a band of 12 hours plus one phase. Reported only, with [EP] tier: progressive reopening of 7 percent at 40 hours, 27 percent at 48 hours and 68 percent at 72 hours. The reopening is not a failure criterion.
\par\smallskip\noindent\textbf{What the current chain does.} Not meaningful yet: depends on Step 3 (post-meal refractory window).
\par\smallskip\noindent\textbf{Checked today by the test suite.}
\begin{itemize}
\item Seeding fidelity: the adult population obeys the conservation identity on every day (adults today = adults yesterday - deaths + emerged, starting from the seeded count), so no seeded individual is dropped or invented; the cohort log and the raster output agree on the number of adults
\item The run completes through the manifest route ('validate\_completeness', resolve, binary) with exit code 0
\end{itemize}
\par\noindent\textbf{What the GIF must show.} One frame per phase, with the time since the meal shown on the female marker and a shaded window of the first 12 hours. A correct run shows no bite or attempt inside the shaded window and the first attempt just after it.

\subsection*{E12 --- Ovarian arrest in the dry season}
\begin{tabular}{@{}ll@{}}
\textit{Mode} & gate \\
\textit{Seed / days} & 20261117 / 30 \\
\textit{Mosquitoes seeded} & 20 female, 0 male \\
\textit{People} & 5; distance from the first mosquito: 141 m, 424 m, 707 m, 728 m, 728 m \\
\textit{Closes} & Step 6 \\
\textit{Runs a behaviour test today} & seeding-only \\
\end{tabular}

\par\smallskip\noindent\textbf{Purpose.} Checks that the fraction of fed females that stop ovarian development during a dry spell (arrest, not death) falls inside the published band for coluzzii. The band is the published data, not a parameter of the scenario.
\par\smallskip\noindent\textbf{Setup.} A group of host-seeking females in one cell, with a few persons nearby, run for thirty days. The environment has a burst of rain followed by a long dry spell, set in the climate input rather than in the scenario file. The cell has no eligible breeding site during the dry spell, so fed females cannot lay.
\par\smallskip\noindent\textbf{What passing or failing implies.} If too few females arrest, dry-season egg production and population are too high and the seasonal dip the data show is missing. If too many arrest, the recovery after the dry season is overstated. Because the band is published data, a fail points to the model, not to a tuned parameter.
\par\smallskip\noindent\textbf{Expected behaviour.} Gate. The fraction of fed females that do not pass stage 2 lies within 19.23 to 34.85 percent, the published range for coluzzii [P]. The scenario measures the fraction and compares it with the band; the band is not adjusted.
\par\smallskip\noindent\textbf{What the current chain does.} Not meaningful yet: depends on Step 6 (arrested state and the dry-spell trigger).
\par\smallskip\noindent\textbf{Checked today by the test suite.}
\begin{itemize}
\item Seeding fidelity: the adult population obeys the conservation identity on every day (adults today = adults yesterday - deaths + emerged, starting from the seeded count), so no seeded individual is dropped or invented; the cohort log and the raster output agree on the number of adults
\item The run completes through the manifest route ('validate\_completeness', resolve, binary) with exit code 0
\end{itemize}
\par\noindent\textbf{What the GIF must show.} One frame per day, with the females coloured by ovarian state (pre-gravid, gravid or arrested) and the rain burst and dry spell drawn as a background band. A correct run shows the arrested share rising during the dry spell and the arrested females still alive when the rain returns.

\subsection*{E13 --- Males and mating supply at one swarm site}
\begin{tabular}{@{}ll@{}}
\textit{Mode} & report \\
\textit{Seed / days} & 20261117 / 3 \\
\textit{Mosquitoes seeded} & 100 female, 100 male \\
\textit{Closes} & none \\
\textit{Runs a behaviour test today} & seeding-only \\
\end{tabular}

\par\smallskip\noindent\textbf{Purpose.} Reports the mating supply one swarm site can offer and how many females remain uninseminated after three nights. It is a report, not a gate, because the plan derives its band by hand from the supply formula before the code is written.
\par\smallskip\noindent\textbf{Setup.} Equal numbers of freshly emerged females and males in a single 1 square km cell, all of whom can reach one swarm site, so every male and female is assigned to it. Some males may stay loose, and the scenario measures how many. The swarm marker density is meant to be fixed in the scenario; the JSON does not yet carry it, and the value enters with Step 8.
\par\smallskip\noindent\textbf{What passing or failing implies.} If the supply is far from the formula, the mating model is inconsistent with the published rates and the uninseminated share is wrong, which matters for population growth after a dry spell. The report lets the team check the number before accepting the swarm model. Because the male count sets the number of pairs, a mismatch appears here first.
\par\smallskip\noindent\textbf{Expected behaviour.} Report band from the formula in section 8.1. Supply per site per night equals the male count times 'p\_mate\_male\_fraction' = 0.072 [P]; for 100 males that is 7.2 pairs. Females present equal the female count times 'female\_attendance\_rate' = 0.1033 [P]; for 100 females that is 10.33. The probability that a female is inseminated is 0.1033 x (1 - exp(-7.2/10.33)) = 0.0518, so about 85.2 percent of females remain uninseminated after three nights [EP]. That figure is a lower bound, because available males decrease as they mate. The earlier published band of 22.3 to 52.8 percent is withdrawn because it cannot be derived from the formula; the implementer computes the report tolerance from this arithmetic.
\par\smallskip\noindent\textbf{What the current chain does.} Not meaningful yet: depends on Step 8 (male site assignment and swarm supply).
\par\smallskip\noindent\textbf{Checked today by the test suite.}
\begin{itemize}
\item Seeding fidelity: the adult population obeys the conservation identity on every day (adults today = adults yesterday - deaths + emerged, starting from the seeded count), so no seeded individual is dropped or invented; the cohort log and the raster output agree on the number of adults
\item The run completes through the manifest route ('validate\_completeness', resolve, binary) with exit code 0
\item Seed sensitivity: two master seeds give different mortality among the 200 seeded adults, so the sub-streams really derive from the single seed
\end{itemize}
\par\noindent\textbf{What the GIF must show.} One frame per night, with the swarm site marked, the males and females near it coloured by assignment and by insemination state, and a running count of inseminated females at the top. A correct run shows the inseminated count rising and levelling off, with a visible share still uninseminated after three nights.

\subsection*{E14 --- Males present before the swarm window}
\begin{tabular}{@{}ll@{}}
\textit{Mode} & gate \\
\textit{Seed / days} & 20261117 / 1 \\
\textit{Mosquitoes seeded} & 1 female, 20 male \\
\textit{Closes} & Step 8 \\
\textit{Runs a behaviour test today} & seeding-only \\
\end{tabular}

\par\smallskip\noindent\textbf{Purpose.} Checks that freshly emerged males assigned to a site do not enter the swarm before their first swarm window, so they cannot inseminate a female early, and that the alarm for uninseminated females is raised.
\par\smallskip\noindent\textbf{Setup.} One host-seeking female with a group of freshly emerged males from the same cell, clustered close to her, run for one day. The day ends before the males' first swarm window. The males carry no state beyond their sex and age. The female is placed close to them, so she is within reach of the site if the site is near.
\par\smallskip\noindent\textbf{What passing or failing implies.} If males enter the swarm before their window, mating happens before the males are physiologically ready and early insemination is overstated. If the alarm is missing, a broken swarm field goes unnoticed, and the alarm is the signal the plan uses to detect that.
\par\smallskip\noindent\textbf{Expected behaviour.} Gate. No female is inseminated by males before their first swarm window. The alarm 'hembras\_sin\_inseminar' is emitted for the cell and night. The female attends the site only if the site lies within her reach.
\par\smallskip\noindent\textbf{What the current chain does.} Not meaningful yet: depends on Step 8 (male site assignment, swarm windows and the alarm).
\par\smallskip\noindent\textbf{Checked today by the test suite.}
\begin{itemize}
\item Seeding fidelity: the adult population obeys the conservation identity on every day (adults today = adults yesterday - deaths + emerged, starting from the seeded count), so no seeded individual is dropped or invented; the cohort log and the raster output agree on the number of adults
\item The run completes through the manifest route ('validate\_completeness', resolve, binary) with exit code 0
\end{itemize}
\par\noindent\textbf{What the GIF must show.} One frame per phase, with the swarm window shaded in the evening. The males are grey before the window and turn to swarm colour only inside it. A correct run shows no male in swarm colour before the shaded window and no insemination during the run.

\subsection*{E15 --- Swarm killing empties the site}
\begin{tabular}{@{}ll@{}}
\textit{Mode} & gate \\
\textit{Seed / days} & 20261117 / 3 \\
\textit{Mosquitoes seeded} & 50 female, 50 male \\
\textit{Closes} & Step 8 \\
\textit{Runs a behaviour test today} & seeding-only \\
\end{tabular}

\par\smallskip\noindent\textbf{Purpose.} Checks that the swarm field is a static, discrete layer: killing the males assigned to one marker removes that site's supply, so the share of uninseminated females rises. It is the test that forces markers to be discrete and removable.
\par\smallskip\noindent\textbf{Setup.} A field of swarm markers with one swarm in the village, equal numbers of males and females, run for three days. The kill of the males assigned to one marker is an intervention the scenario applies. The JSON does not express the kill yet; the flag enters with Step 8.
\par\smallskip\noindent\textbf{What passing or failing implies.} If killing the males at a marker does not reduce supply, the swarm is a number rather than a place, and the model could not represent a targeted intervention. A pass supports the static site layer; a fail means site assignment does not reach the supply.
\par\smallskip\noindent\textbf{Expected behaviour.} Gate, directional. After the kill the supply of the killed site falls to zero, and the share of uninseminated females is higher than in the same-seed run without the kill. The plan gives no numeric threshold for the rise, so the test checks the direction of the change and that the killed site's supply is removed.
\par\smallskip\noindent\textbf{What the current chain does.} Not meaningful yet: depends on Step 8 (site supply and the swarm-killing flag).
\par\smallskip\noindent\textbf{Checked today by the test suite.}
\begin{itemize}
\item Seeding fidelity: the adult population obeys the conservation identity on every day (adults today = adults yesterday - deaths + emerged, starting from the seeded count), so no seeded individual is dropped or invented; the cohort log and the raster output agree on the number of adults
\item The run completes through the manifest route ('validate\_completeness', resolve, binary) with exit code 0
\item Seed sensitivity: two master seeds give different mortality among the seeded adults
\end{itemize}
\par\noindent\textbf{What the GIF must show.} One frame per night, with the markers as fixed points, the males and females coloured by assignment, and the killed marker crossed out after the kill. A correct run shows the killed site's males disappearing and its females staying uninseminated while other sites keep mating.

\subsection*{E16 --- Nets at 60 percent do not offset swarm supply}
\begin{tabular}{@{}ll@{}}
\textit{Mode} & report \\
\textit{Seed / days} & 20261117 / 3 \\
\textit{Mosquitoes seeded} & 50 female, 50 male \\
\textit{Closes} & none \\
\textit{Runs a behaviour test today} & seeding-only \\
\end{tabular}

\par\smallskip\noindent\textbf{Purpose.} Counter-scenario that records a model limit: individual protection from insecticide-treated nets at 60 percent coverage does not restore the swarm supply lost at a site. It is a report because the plan treats the result as a documented limit, not as a model output to tune.
\par\smallskip\noindent\textbf{Setup.} The same swarm-site field as the swarm-killing scenario, with equal numbers of males and females, run for three days, with net coverage set to 60 percent. The net coverage belongs to the scenario definition; the JSON does not carry a net field yet.
\par\smallskip\noindent\textbf{What passing or failing implies.} If nets appeared to restore supply, the model would be suggesting that individual protection can compensate for lost swarm sites, which the plan does not support. Keeping the limit in the report stops a reader from taking net coverage as a mating-supply intervention.
\par\smallskip\noindent\textbf{Expected behaviour.} Report only, with no band in the plan. The report records the uninseminated share and the supply per site with nets and without, and states that the nets do not compensate for the supply loss. The plan records this as a documented limit of the model, not as a result.
\par\smallskip\noindent\textbf{What the current chain does.} Not meaningful yet: depends on Step 8 (the swarm supply the nets are compared against).
\par\smallskip\noindent\textbf{Checked today by the test suite.}
\begin{itemize}
\item Seeding fidelity: the adult population obeys the conservation identity on every day (adults today = adults yesterday - deaths + emerged, starting from the seeded count), so no seeded individual is dropped or invented; the cohort log and the raster output agree on the number of adults
\item The run completes through the manifest route ('validate\_completeness', resolve, binary) with exit code 0
\end{itemize}
\par\noindent\textbf{What the GIF must show.} Same layout as the swarm-killing GIF, with a net coverage indicator at 60 percent drawn on the frame. A correct run shows the sites that lose supply staying weak with nets in place.

\subsection*{E17 --- Sterile males compete for site supply}
\begin{tabular}{@{}ll@{}}
\textit{Mode} & gate \\
\textit{Seed / days} & 20261117 / 3 \\
\textit{Mosquitoes seeded} & 50 female, 50 male \\
\textit{Closes} & Step 10 \\
\textit{Runs a behaviour test today} & seeding-only \\
\end{tabular}

\par\smallskip\noindent\textbf{Purpose.} Checks that sterile males assigned to the same swarm sites as the fertile males take a proportional share of the supply, so that supply per site falls with their number and the uninseminated share rises. It is the sterile-insect test on site supply.
\par\smallskip\noindent\textbf{Setup.} Equal numbers of females and males assigned to the swarm sites, with sterile males added and assigned to the same sites, so they compete for the same supply. Run for three days. The sterility flag is not in the JSON yet; it enters with Step 10.
\par\smallskip\noindent\textbf{What passing or failing implies.} If supply does not fall with the number of sterile males assigned, the model cannot represent sterile-insect release, an intervention a programme may run. Competition among individual sterile males inside a swarm is not simulated (N11, section 8.6), so the test covers supply per site, not competition within the swarm.
\par\smallskip\noindent\textbf{Expected behaviour.} Gate. Supply per site falls in proportion to the number of sterile males assigned to it, and the uninseminated female share rises. The Step 10 criterion also requires that two identical runs match byte for byte (Level 5).
\par\smallskip\noindent\textbf{What the current chain does.} Not meaningful yet: depends on Step 10 (sterile males on the site supply).
\par\smallskip\noindent\textbf{Checked today by the test suite.}
\begin{itemize}
\item Seeding fidelity: the adult population obeys the conservation identity on every day (adults today = adults yesterday - deaths + emerged, starting from the seeded count), so no seeded individual is dropped or invented; the cohort log and the raster output agree on the number of adults
\item The run completes through the manifest route ('validate\_completeness', resolve, binary) with exit code 0
\end{itemize}
\par\noindent\textbf{What the GIF must show.} One frame per night, with fertile and sterile males drawn in different colours at each site and the supply of each site shown as a bar. A correct run shows each site's supply shrinking as its sterile share grows, and the uninseminated count rising with it.

\subsection*{E18 --- Wind deflection and night blocking}
\begin{tabular}{@{}ll@{}}
\textit{Mode} & gate \\
\textit{Seed / days} & 20261117 / 3 \\
\textit{Mosquitoes seeded} & 1 female, 0 male \\
\textit{Closes} & Step 5 \\
\textit{Runs a behaviour test today} & seeding-only \\
\end{tabular}

\par\smallskip\noindent\textbf{Purpose.} Checks the wind branch of the movement law: a female flying against the wind is deflected upwind, and forward progress is blocked when the tail-wind component reaches her flight speed. It is the gate for the wind branch.
\par\smallskip\noindent\textbf{Setup.} One host-seeking female at the centre of the synthetic study area, run for three days. The synthetic area carries a wind field of the 100 m wind components, loaded at run time. The wind is not written in the JSON; it comes from the synthetic area inputs.
\par\smallskip\noindent\textbf{What passing or failing implies.} If the female ignores the wind, the model cannot show how wind affects dispersal and therefore the reach of transmission. If the block never occurs, females would fly against strong winds at full speed. The wind is taken at 100 m, not at flight height, so the branch can only be validated against the synthetic area; this is a declared limit.
\par\smallskip\noindent\textbf{Expected behaviour.} Gate. The flight is deflected upwind by up to 30 degrees ('anemotaxis\_rule', [EP]). Forward progress is blocked when the tail-wind component, projected on the heading, reaches the flight speed of 0.225 to 0.45 m/s [P]. The threshold is computed at run time from the configured flight speed and is never hard-coded. A permanent night block under the measured wind is accepted as a documented result of the 100 m level, not as an error.
\par\smallskip\noindent\textbf{What the current chain does.} Not meaningful yet: depends on Step 5 (wind branch of the cascade). Without a loaded wind field the wind rule is inactive.
\par\smallskip\noindent\textbf{Checked today by the test suite.}
\begin{itemize}
\item Seeding fidelity: the adult population obeys the conservation identity on every day (adults today = adults yesterday - deaths + emerged, starting from the seeded count), so no seeded individual is dropped or invented; the cohort log and the raster output agree on the number of adults
\item The run completes through the manifest route ('validate\_completeness', resolve, binary) with exit code 0
\end{itemize}
\par\noindent\textbf{What the GIF must show.} One frame per phase, with the female track and wind arrows drawn at each frame, the arrow length showing wind speed. A correct run shows the track bent upwind relative to the wind arrows, and forward progress stopping when the tail-wind arrow reaches the flight speed.
